% GENERATED FILE. Edit canonical lessons, then run npm run generate:books.
% canonical-source-hash: fnv1a64:b42cf186753eadb8
% canonical-lessons: AR-C395-marji, AR-C395-istashhada, AR-C395-iqtabasa, AR-C395-talaa, AR-C395-taammala

\chapter{Reading and quoting}
\label{ch:ar-reading-and-quoting}
\hlchaptermodality{395}

\begin{chapteropening}
\textbf{By the end of this chapter:} \emph{I can talk about a reference, and say that someone cited, quoted, read up on or reflected.}

Complete the last lesson of chapter 395: I can talk about a reference, and say that someone cited, quoted, read up on or reflected.
\end{chapteropening}

\section[\texorpdfstring{\ar{مرجع} (marjiʿ)}{مرجع (marjiʿ)}]{\ar{مرجع} (marjiʿ) — a reference}

\label{lesson:AR-C395-marji}



\begin{quote}
Before the new one: say the Arabic for to summarize, then the Arabic for a summary.
\end{quote}

\subsection*{You'll want to know: \ar{مرجع}}
\textbf{\ar{مرجع}} — \emph{marjiʿ} — \textquotedblleft{}a reference\textquotedblright{}.

\begin{quote}
\ar{هذا} \ar{القاموس} \ar{مرجع} \ar{مهم}.
\end{quote}

\begin{quote}
\emph{hādhā al-qāmūs marjiʿ muhimm.}
\end{quote}

\begin{quote}
This dictionary is an important reference.
\end{quote}

Plural: \ar{مراجع}, \emph{marājiʿ}.

\begin{cousinweb}
\textbf{\ar{مرجع}} (\emph{marjiʿ}) is built on the native Arabic root \emph{r-j-ʿ}, \textquotedblleft{}to return, to go back\textquotedblright{} (Lane). \emph{Marjiʿ} is a noun of place, a place one goes back to: a reference.
\end{cousinweb}

\begin{grammarlens}[title={What you've built}]
One more word for giving reasons and opinions.
\end{grammarlens}

\subsection*{Your turn}
\begin{itemize}
\raggedright
  \item \emph{Say it:} \emph{marjiʿ}
  \item \emph{Say it:} \emph{marjiʿ}, once more
  \item \emph{Say it:} \emph{hādhā al-qāmūs marjiʿ muhimm.} — This dictionary is an important reference.
  \item \emph{Recall:} say the Arabic for to summarize, then the Arabic for a summary, then say \emph{marjiʿ} again
\end{itemize}

\begin{tcolorbox}[breakable,colback=teal!4,colframe=teal!35!black,title={Before you move on}]
What does \ar{مرجع} mean? (\textquotedblleft{}A reference\textquotedblright{}.) Say it once more.
\end{tcolorbox}
\section[\texorpdfstring{\ar{استشهد} (istashhada)}{استشهد (istashhada)}]{\ar{استشهد} (istashhada) — to cite}

\label{lesson:AR-C395-istashhada}



\begin{quote}
Before the new one: say the Arabic for a summary, then the Arabic for a reference.
\end{quote}

\subsection*{You'll want to know: \ar{استشهد}}
\textbf{\ar{استشهد}} — \emph{istashhada} — \textquotedblleft{}to cite\textquotedblright{}.

\begin{quote}
\ar{استشهد} \ar{الكاتب} \ar{بكلام} \ar{جده}.
\end{quote}

\begin{quote}
\emph{istashhada al-kātib bi-kalām jaddih.}
\end{quote}

\begin{quote}
The writer cited his grandfather's words.
\end{quote}

Present: \ar{يستشهد}, \emph{yastashhid}.

It takes \ar{بـ}: \ar{استشهد} \ar{بـ} …. In another sense it means to die as a martyr.

\begin{cousinweb}
\textbf{\ar{استشهد}} (\emph{istashhada}) is built on the native Arabic root \emph{sh-h-d}, \textquotedblleft{}to witness, to be present\textquotedblright{} (Lane). It is Form X, to call for a witness: to cite as evidence.
\end{cousinweb}

\begin{grammarlens}[title={What you've built}]
One more word for giving reasons and opinions.
\end{grammarlens}

\subsection*{Your turn}
\begin{itemize}
\raggedright
  \item \emph{Say it:} \emph{istashhada}
  \item \emph{Say it:} \emph{istashhada}, once more
  \item \emph{Say it:} \emph{istashhada al-kātib bi-kalām jaddih.} — The writer cited his grandfather's words.
  \item \emph{Recall:} say the Arabic for a summary, then the Arabic for a reference, then say \emph{istashhada} again
\end{itemize}

\begin{tcolorbox}[breakable,colback=teal!4,colframe=teal!35!black,title={Before you move on}]
What does \ar{استشهد} mean? (\textquotedblleft{}To cite\textquotedblright{}.) Say it once more.
\end{tcolorbox}
\section[\texorpdfstring{\ar{اقتبس} (iqtabasa)}{اقتبس (iqtabasa)}]{\ar{اقتبس} (iqtabasa) — to quote}

\label{lesson:AR-C395-iqtabasa}



\begin{quote}
Before the new one: say the Arabic for a reference, then the Arabic for to cite.
\end{quote}

\subsection*{You'll want to know: \ar{اقتبس}}
\textbf{\ar{اقتبس}} — \emph{iqtabasa} — \textquotedblleft{}to quote\textquotedblright{}.

\begin{quote}
\ar{اقتبس} \ar{الصحفي} \ar{جملة} \ar{من} \ar{الخطاب}.
\end{quote}

\begin{quote}
\emph{iqtabasa aṣ-ṣaḥafī jumla min al-khiṭāb.}
\end{quote}

\begin{quote}
The journalist quoted a sentence from the speech.
\end{quote}

Present: \ar{يقتبس}, \emph{yaqtabis}.

\begin{cousinweb}
\textbf{\ar{اقتبس}} (\emph{iqtabasa}) is built on the native Arabic root \emph{q-b-s}, \textquotedblleft{}to take fire from\textquotedblright{} (Lane). Lane gives the root for taking a flame from a fire, and also for taking knowledge: Form VIII came to mean to quote.
\end{cousinweb}

\begin{grammarlens}[title={What you've built}]
One more word for giving reasons and opinions.
\end{grammarlens}

\subsection*{Your turn}
\begin{itemize}
\raggedright
  \item \emph{Say it:} \emph{iqtabasa}
  \item \emph{Say it:} \emph{iqtabasa}, once more
  \item \emph{Say it:} \emph{iqtabasa aṣ-ṣaḥafī jumla min al-khiṭāb.} — The journalist quoted a sentence from the speech.
  \item \emph{Recall:} say the Arabic for a reference, then the Arabic for to cite, then say \emph{iqtabasa} again
\end{itemize}

\begin{tcolorbox}[breakable,colback=teal!4,colframe=teal!35!black,title={Before you move on}]
What does \ar{اقتبس} mean? (\textquotedblleft{}To quote\textquotedblright{}.) Say it once more.
\end{tcolorbox}
\section[\texorpdfstring{\ar{طالع} (ṭālaʿa)}{طالع (ṭālaʿa)}]{\ar{طالع} (ṭālaʿa) — to read up on}

\label{lesson:AR-C395-talaa}



\begin{quote}
Before the new one: say the Arabic for to cite, then the Arabic for to quote.
\end{quote}

\subsection*{You'll want to know: \ar{طالع}}
\textbf{\ar{طالع}} — \emph{ṭālaʿa} — \textquotedblleft{}to read up on\textquotedblright{}.

\begin{quote}
\ar{يطالع} \ar{أبي} \ar{الجريدة} \ar{كل} \ar{صباح}.
\end{quote}

\begin{quote}
\emph{yuṭāliʿu abī al-jarīda kulla ṣabāḥ.}
\end{quote}

\begin{quote}
My father reads through the newspaper every morning.
\end{quote}

Present: \ar{يطالع}, \emph{yuṭāliʿ}.

\begin{cousinweb}
\textbf{\ar{طالع}} (\emph{ṭālaʿa}) is built on the native Arabic root \emph{ṭ-l-ʿ}, \textquotedblleft{}to rise (of the sun)\textquotedblright{} (Lane). It is Form III, to look into a thing closely: to read up on, to peruse.
\end{cousinweb}

\begin{grammarlens}[title={What you've built}]
One more word for giving reasons and opinions.
\end{grammarlens}

\subsection*{Your turn}
\begin{itemize}
\raggedright
  \item \emph{Say it:} \emph{ṭālaʿa}
  \item \emph{Say it:} \emph{ṭālaʿa}, once more
  \item \emph{Say it:} \emph{yuṭāliʿu abī al-jarīda kulla ṣabāḥ.} — My father reads through the newspaper every morning.
  \item \emph{Recall:} say the Arabic for to cite, then the Arabic for to quote, then say \emph{ṭālaʿa} again
\end{itemize}

\begin{tcolorbox}[breakable,colback=teal!4,colframe=teal!35!black,title={Before you move on}]
What does \ar{طالع} mean? (\textquotedblleft{}To read up on\textquotedblright{}.) Say it once more.
\end{tcolorbox}
\section[\texorpdfstring{\ar{تأمل} (taʾammala)}{تأمل (taʾammala)}]{\ar{تأمل} (taʾammala) — to reflect}

\label{lesson:AR-C395-taammala}



\begin{quote}
Before the new one: say the Arabic for to quote, then the Arabic for to read up on.
\end{quote}

\subsection*{You'll want to know: \ar{تأمل}}
\textbf{\ar{تأمل}} — \emph{taʾammala} — \textquotedblleft{}to reflect\textquotedblright{}.

\begin{quote}
\ar{تأمل} \ar{الرجل} \ar{البحر} \ar{طويلا}.
\end{quote}

\begin{quote}
\emph{taʾammala ar-rajul al-baḥr ṭawīlan.}
\end{quote}

\begin{quote}
The man contemplated the sea for a long time.
\end{quote}

Present: \ar{يتأمل}, \emph{yataʾammal}.

It takes a direct object or \ar{في}: \ar{تأمل} \ar{في} ….

\begin{cousinweb}
\textbf{\ar{تأمل}} (\emph{taʾammala}) is built on the native Arabic root \emph{ʾ-m-l}, \textquotedblleft{}to hope for\textquotedblright{} (Lane). It is Form V: to gaze on a thing long and carefully, to reflect on it.
\end{cousinweb}

\begin{grammarlens}[title={What you've built}]
One more word for giving reasons and opinions.
\end{grammarlens}

\subsection*{Your turn}
\begin{itemize}
\raggedright
  \item \emph{Say it:} \emph{taʾammala}
  \item \emph{Say it:} \emph{taʾammala}, once more
  \item \emph{Say it:} \emph{taʾammala ar-rajul al-baḥr ṭawīlan.} — The man contemplated the sea for a long time.
  \item \emph{Recall:} say the Arabic for to quote, then the Arabic for to read up on, then say \emph{taʾammala} again
\end{itemize}

\begin{tcolorbox}[breakable,colback=teal!4,colframe=teal!35!black,title={Before you move on}]
What does \ar{تأمل} mean? (\textquotedblleft{}To reflect\textquotedblright{}.) Say it once more.
\end{tcolorbox}
